\section{Discussion}
\label{sec:discussion}

\subsection{Why Is Decay Instantaneous?}

Our results show that the model completely absorbs a steering perturbation within
a single forward pass. We attribute this to two architectural features.

First, \textbf{LayerNorm re-normalizes activations} at each layer. When a
steering vector shifts the hidden state's direction and magnitude, subsequent
LayerNorm operations project the result back toward the learned activation
distribution. \citet{dang2026selective} identify this norm-distortion mechanism
and show that it compounds across layers. Our results confirm that by the time
the perturbed hidden state has passed through the remaining layers above the
injection point, the normalization has completely absorbed the perturbation.

Second, \textbf{residual connections and the learned weight matrices} implement
a form of manifold projection. \citet{xu2026why} propose that model activations
lie on a low-dimensional manifold and that subsequent layers project perturbed
activations back onto this manifold. Our finding that decay is instantaneous
regardless of perturbation magnitude (multiplier ablation) or duration (duration
ablation) is consistent with this view: the projection is strong enough to
absorb even $5\times$ perturbations in a single step.

\subsection{Representational vs.\ Behavioral Persistence}

The dissociation between representational decay and behavioral persistence is
the most practically significant finding. Hidden states forget the steering
vector immediately, but the tokens generated under steering remain in the
context window and continue to influence the model's output distribution. This
behavioral persistence has a half-life exceeding 75~tokens (\figref{fig:behavioral})---two orders of
magnitude longer than the representational half-life of less than one token.

This finding aligns with the Bayesian framework of \citet{bigelow2025belief}:
the steering vector shifts the model's prior, which then gets ``baked into'' the
generated tokens. Once those tokens are in the context, they provide ongoing
evidence for the steered behavior, sustaining it even after the representational
prior shift has been completely absorbed. The \TFclean condition confirms this
interpretation: when unsteered tokens replace the steered tokens in context,
all behavioral effects disappear.

\subsection{Implications for Practitioners}

\para{Continuous application is necessary, not just robust.}
Our results explain why standard methods like \caa apply steering vectors at
every token position. This is not merely a conservative design choice---it is
necessary because a single-application perturbation is fully absorbed within one
step. Practitioners cannot reduce computational cost by applying vectors
intermittently.

\para{Brief steering can have lasting behavioral effects.}
Despite the instantaneous representational decay, even brief steering ($K = 1$
to $3$ tokens) at the start of generation can create lasting behavioral effects
through the generated token sequence. If the goal is to set a behavioral
trajectory rather than maintain continuous representational alignment, short
initial steering may suffice.

\para{The token sequence is the primary control surface.}
Our results suggest that for long-term behavioral control, the generated token
sequence---not hidden state manipulation---is the most durable mechanism.
Methods that operate through the context window (\eg prompt engineering, few-shot
examples) may produce more persistent effects than activation-level
interventions.

\subsection{Implications for Theory}

Our findings inform several theoretical questions about steering:

\para{Steering as dynamic bias.}
\citet{xu2026why} characterize steering as adding a dynamic bias to model
weights. Our results show that this bias has no temporal memory: it affects
only the forward pass in which it is applied. Each forward pass is
independent---the model does not accumulate or propagate the bias across steps.

\para{The PID framework.}
\citet{nguyen2025pid} frame standard steering as P-control with inherent
steady-state error. Our results sharpen this analysis: the ``steady-state error''
of a removed P-controller is 100\%. The integral term in PID steering is not
merely reducing error---it is providing the \emph{entire} persistent signal,
since the proportional term contributes nothing once removed.

\para{Transformer robustness.}
The instantaneous absorption of perturbations---even at $5\times$ the default
multiplier---suggests that transformer hidden states are remarkably robust. The
model's learned dynamics actively resist deviation from the natural activation
manifold, which may have implications beyond steering for understanding
adversarial robustness and representation stability.

\subsection{Limitations}

\para{Single model.}
We test only \gptxl. Larger models with more layers, different normalization
schemes (RMSNorm vs.\ LayerNorm), or different architectures may show different
decay dynamics. In particular, models with more layers may show slower decay
if the manifold projection is weaker.

\para{Single injection layer.}
We inject at one layer per experiment. Multi-layer steering, as used in standard
\caa, applies vectors at multiple layers simultaneously and may produce
different decay patterns due to reinforcement across layers.

\para{Greedy decoding.}
We use argmax decoding throughout. Sampling-based decoding introduces
stochasticity that could interact with decay dynamics, potentially amplifying
or dampening the effect of steered tokens in context.

\para{Hidden state measurement at injection layer only.}
We measure cosine similarity at the injection layer. The perturbation may
propagate to other layers with different dynamics, though the output-level KL
divergence results (which reflect all-layer effects) are consistent with
instantaneous decay.

\para{Small contrastive sets.}
We use 20 contrastive pairs per behavior, which is on the lower end for steering
vector extraction. Larger sets could produce more robust vectors with
potentially different decay properties, though we see consistent results across
all three behaviors tested.
